\let\negmedspace\undefined
\let\negthickspace\undefined
\documentclass[journal]{article}
\usepackage[a5paper, margin=10mm, onecolumn]{geometry}
\usepackage{lmodern} % Ensure lmodern is loaded for pdflatex

\setlength{\headheight}{1cm} % Set the height of the header box
\setlength{\headsep}{0mm}     % Set the distance between the header box and the top of the text

\usepackage{gvv-book}
\usepackage{gvv}
\usepackage{cite}
\usepackage{textcomp}
\usepackage{amsmath,amssymb,amsfonts,amsthm}
\usepackage{algorithmic}
\usepackage{graphicx}
\graphicspath{{./figs/}}
\usepackage{textcomp}
\usepackage{xcolor}
\usepackage{txfonts}
\usepackage{listings}
\usepackage{enumitem}
\usepackage{mathtools}
\usepackage{gensymb}
\usepackage{comment}
\usepackage[breaklinks=true]{hyperref}
\usepackage{tkz-euclide} 
\usepackage{listings}
\usepackage{gvv}                                        
\def\inputGnumericTable{}                                 
\usepackage[latin1]{inputenc}                                
\usepackage{color}                                            
\usepackage{array}                                            
\usepackage{longtable}                                       
\usepackage{calc}                                             
\usepackage{multirow}                                         
\usepackage{hhline}                                           
\usepackage{ifthen}                                           
\usepackage{lscape}
\usepackage{circuitikz}
\tikzstyle{block} = [rectangle, draw, fill=blue!20, 
text width=4em, text centered, rounded corners, minimum height=3em]
\tikzstyle{sum} = [draw, fill=blue!10, circle, minimum size=1cm, node distance=1.5cm]
\tikzstyle{input} = [coordinate]
\tikzstyle{output} = [coordinate]


\begin{document}
	
	\bibliographystyle{IEEEtran}
	\vspace{3cm}
	
\title{10.6.2}
\author{EE25BTECH11047 - RAVULA SHASHANK REDDY}
\maketitle
\hrulefill
\bigskip 

\renewcommand{\thetable}{\theenumi}
\setlength{\intextsep}{10pt}

\textbf{Question:} \\

Construct a pair of tangents to a circle of radius \(4\,cm\) from a point \(P\) at a distance of 

\(7\,cm\) from its centre \(O\).

\textbf{Solution:}\\


The tangent directions \(\vec{m}\) from an external point \(\vec{h}\) to the circle 
\begin{align}
g(\vec{x}) = \vec{x}^\top \vec{x} - r^2 = 0
\end{align}

satisfy 
\begin{align}
\vec{m}^\top \vec{\Phi} \vec{m} = 0
\end{align}

where 
\begin{align}
\vec{\Phi} &= \vec{h}\vec{h}^\top - g(\vec{h})\vec{I} \label{eq:sigma_def}
\end{align}

With the point \(\vec{h} = d\vec{e}_1\), we have \(g(\vec{h}) = d^2 - r^2\). From \eqref{eq:sigma_def},
\begin{align}
\vec{\Phi} &= 
\myvec{
d^2 & 0 \\
0 & 0
}
-
(d^2 - r^2)
\myvec{
1 & 0 \\
0 & 1
}
\\
&=
\myvec{
r^2 & 0 \\
0 & -(d^2 - r^2)
} \label{eq:sigma_matrix}
\end{align}

Since \(\vec{\Phi}\) is a diagonal matrix, its eigenvalues are:
\begin{align}
\lambda_1 &= r^2, \quad \lambda_2 = -(d^2 - r^2) \label{eq:eigenvalues}
\end{align}

The matrix of orthonormal eigenvectors is 
\[
\vec{P} = \vec{I}
\]

Hence, the unit direction vectors of the tangents are given by:
\begin{align}
\vec{m} = \frac{1}{\sqrt{\lambda_1 - \lambda_2}} \, \vec{P} 
\myvec{
\sqrt{-\lambda_2} \\[4pt]
\pm \sqrt{\lambda_1}
}
\label{eq:m_general}
\end{align}

From \eqref{eq:eigenvalues},
\begin{align}
\lambda_1 - \lambda_2 &= r^2 - (-(d^2 - r^2)) = d^2 \\
-\lambda_2 &= d^2 - r^2
\end{align}

Substituting into \eqref{eq:m_general}:
\begin{align}
\vec{m} 
&= \frac{1}{\sqrt{d^2}} 
\myvec{
\sqrt{d^2 - r^2} \\[4pt]
\pm r
}
\\
&= \frac{1}{d}
\myvec{
\sqrt{d^2 - r^2} \\[4pt]
\pm r
}
\label{eq:m_final}
\end{align}

The points of contact are given by 
\begin{align}
\vec{q} = \vec{h} + \kappa \vec{m}
\end{align}

Satisfies the circle equation
\begin{align}
g(\vec{h} + \kappa\vec{m}) 
&= (\vec{h} + \kappa\vec{m})^\top (\vec{h} + \kappa\vec{m}) - r^2 = 0 \\
&\implies 
\kappa^2 (\vec{m}^\top \vec{m}) + 2\kappa(\vec{h}^\top \vec{m}) + g(\vec{h}) = 0
\end{align}

For a tangent, \(\kappa\) has a single repeated root.
Since \(\vec{m}\) is a unit vector (\(\vec{m}^\top \vec{m} = 1\)):
\begin{align}
\kappa &= -(\vec{h}^\top \vec{m}) \label{eq:kappa_def}
\end{align}

\begin{align}
\vec{h}^\top \vec{m} 
&= (d\vec{e}_1)^\top 
\frac{1}{d}
\myvec{
\sqrt{d^2 - r^2} \\[4pt]
\pm r
}
= \sqrt{d^2 - r^2} 
\end{align}

Hence,
\[
\kappa = -\sqrt{d^2 - r^2}
\]

The points of contact are:
\begin{align}
\vec{q} 
&= d\vec{e}_1 - \sqrt{d^2 - r^2}
\frac{1}{d}
\myvec{
\sqrt{d^2 - r^2} \\[4pt]
\pm r
}
\\[6pt]
&=
\myvec{
\dfrac{r^2}{d} \\[6pt]
\mp \dfrac{r\sqrt{d^2 - r^2}}{d}
}
\label{eq:points}
\end{align}

\textbf{Substituting \(r = 4, \; d = 7\):}
\begin{align}
\sqrt{d^2 - r^2} &= \sqrt{49 - 16} = \sqrt{33} \\
\vec{q} 
&=
\myvec{
\dfrac{16}{7} \\[6pt]
\mp \dfrac{4\sqrt{33}}{7}
}
\end{align}

Thus, the points of contact of the tangents are:
\[
\boxed{
\vec{q}_1 = 
\myvec{
\dfrac{16}{7} \\[4pt]
\dfrac{4\sqrt{33}}{7}
}, 
\quad
\vec{q}_2 =
\myvec{
\dfrac{16}{7} \\[4pt]
-\dfrac{4\sqrt{33}}{7}
}
}
\]

\newpage
\begin{figure}
    \centering
    \includegraphics[width=1.0\linewidth]{figs/Fig_1.png}
    \caption{}
    \label{fig:placeholder}
\end{figure}
 
\end{document}